% D04-003. Начало изложения по AI-02-RU P0077--P0102.
\section{Представления токов и напряжений на входе модели}
\label{sec:ch4_signal_representations}

\textit{Ниже изложены подготовка сигналов и способы их представления. Количественное сопоставление вариантов ещё в работе и отмечено отдельно в конце раздела.}


\iffalse
(Примечание автора: Как-то тоже ну ОЧЕНЬ урезано, если ты пишешь это как начальный кусок, то пожалуйста, делай так, чтобы и в ПДФ было видно, чтобы я не придавал сильного значения и понимал, что ты ещё будешь много делать. А то прям хочется написать, что ничего тут нет, давай править, картинки вставлять и т.д..) **[AA069: закрыто автором; сохранено в истории]**
(Примечание автора: Замечание закрыть, достаточно фразы о том, что процесс идёт и вычитка не требуется.)
\fi
\authorreviewnote{AA069}


Архитектура определяет способ обработки, а входное представление
определяет, какие величины эта обработка получает.
Исследование признакового пространства\footnote{Евдаков А.Е.,
Яблоков А.А., Филатова Г.А. Формирование оптимального признакового
пространства нейросетей для задач анализа осциллограмм переходных
процессов в электроэнергетических системах.} рассматривает
совместно состав измерений, способ преобразования сигналов
и доступную предысторию. Сопоставляются семейства моделей,
описанные в разделе~\ref{sec:ch4_physkan_architectures}.
Это позволяет разделить изменения самого входа и изменения
вычислительной структуры сети.

Для БАВР выбор представления связан с характером переходного
процесса. Мгновенные значения сохраняют форму фронта и искажения,
а спектральные признаки выделяют амплитудные и фазовые соотношения.
История нескольких периодов позволяет описать развитие режима,
но объём передаваемых данных и время формирования признаков
также входят в оценку алгоритма.

\subsection{Состав каналов и восстановление отсутствующих измерений}
\label{sec:ch4_channel_preparation}

В исследовании признаковых пространств используется общий
набор из восьми сигналов. В него входят фазные токи $i_A$,
$i_B$, $i_C$, ток нулевой последовательности и соответствующие
три фазных напряжения с напряжением нулевой последовательности.
Общий состав нужен для сопоставления записей с разными
наборами зарегистрированных каналов. Сведения о том, какие
величины непосредственно измерены и какие восстановлены,
при этом имеют разный смысл.

(Примечание автора: ниже про отсутствующий канал, проверь, а так ли у нас в коде, действительно ли мы IN учитываем в формуле всегда, или принимаем сеть изолированной и восстанавливает ток фазы В? А может и вообще его зануляем. Вот чёт не помню, и как бы не зануляли его в последних версиях…)

Отсутствующий канал может быть получен расчётом, если известны
связанные с ним измерения и их направления. Например, при
выбранном направлении тока нейтрали $i_N$ баланс токов узла
имеет вид $i_A+i_B+i_C+i_N=0$. Тогда отсутствующий ток фазы B
определяется выражением
\begin{equation}
 i_B(t)=-\bigl(i_A(t)+i_C(t)+i_N(t)\bigr).
 \label{eq:ch4_current_reconstruction}
\end{equation}
Знак нейтрального тока зависит от принятого направления
его измерения. Сигнал $i_N$ и составляющая $i_0$ также
различаются масштабом и, в зависимости от соглашения,
полярностью. Поэтому название канала не заменяет проверки
его физического определения.

Для напряжений различаются фазные и линейные измерения.
Разности фазных напряжений дают линейные, но только по линейным
напряжениям нельзя определить общую нулевую составляющую.
При отсутствии дополнительной информации полное восстановление
фазного набора неоднозначно. Это особенно существенно для ОЗЗ,
при котором смещение нейтрали является частью исследуемого процесса.
Расчётная величина не считается независимым новым измерением
и сохраняет ошибки исходных каналов.

Если необходимых исходных сигналов нет, значение остаётся
отсутствующим. Его нельзя без пояснения приравнивать измеренному
нулю. В моделях с поддержкой пропусков используется отдельный
маркер отсутствия. Влияние состава доступных каналов проверяется
как самостоятельное условие испытания, а не выводится только
из возможности вычислить входной вектор нужной длины.

\subsection{Временное окно и доступная предыстория}
\label{sec:ch4_feature_history}

При передаче мгновенных значений входное окно содержит
последовательные отсчёты всех выбранных каналов. Такое
представление сохраняет несинусоидальные участки, броски тока
и взаимное положение изменений токов и напряжений.
Одновременно модель получает зависимость от частоты
дискретизации, начальной фазы и положения события внутри окна.
Согласование записей поэтому относится не только к числу каналов,
но и к физической шкале времени.

Длина доступного фрагмента и длительность окна вычисления
спектра являются разными параметрами. Из одного фрагмента
можно сформировать последовательность спектральных векторов,
каждый из которых получен по своему локальному окну.
Модель тогда видит изменение признаков во времени,
а не только один усреднённый спектральный срез.
Шаг между такими векторами задаёт частоту обновления входа.
Он не совпадает автоматически с длительностью всего буфера.

(Примечание автора: ну вроде бы всё верно, но можно бы и упростить чуть описание. Может даже к всему подразделу этому)

Расположение окна относительно оцениваемой метки также
существенно. Если правая граница окна совпадает с текущим
моментом, в признаки входят уже полученные отсчёты.
При использовании центрированного окна для той же метки
потребуются более поздние отсчёты. Эти варианты могут давать
разное качество на сохранённых осциллограммах, но имеют разные
условия доступности данных при работе измерительного органа.

Для классификации окна оценивается выбранный фрагмент целиком.
Для зонной детекции выход относится к определённой части
последовательности. Поэтому увеличение предыстории не задаёт
само по себе правило разметки и не определяет задержку
обнаружения события. Эта связь рассматривается совместно
с единицей оценки из раздела~\ref{sec:ch4_evaluation_methodology}.

(Примечание автора: А как это там задаётся, в том разделе? В общем считаю нужным чуть переписать раздел.)

\subsection{Спектральные векторы и выбор координат}
\label{sec:ch4_spectral_coordinates}

В спектральном представлении из окна выделяются гармонические
составляющие сигналов. В исследовании признакового пространства
использованы первые девять гармоник промышленной частоты.
Число доступных составляющих ограничивается частотой
дискретизации исходной записи. Длина окна, весовая функция
и коэффициент нормировки преобразования влияют на получаемые
амплитуды и учитываются при сопоставлении вариантов.

Один комплексный коэффициент можно передать сети двумя
вещественными компонентами. В декартовом представлении это
действительная и мнимая части. В полярном представлении это
модуль и аргумент
\begin{equation}
 \dot{X}=X_{\mathrm{Re}}+jX_{\mathrm{Im}}
        =A\bigl(\cos\varphi+j\sin\varphi\bigr).
 \label{eq:ch4_feature_coordinates}
\end{equation}
Обе формы описывают одну комплексную величину, но задают
разные входные зависимости для обучаемого преобразования.
В декартовой форме изменение фазы перераспределяет значения
между двумя компонентами. В полярной форме модуль и угол
выделены явно и могут обрабатываться раздельными ветвями.

Полярные координаты удобны для введения комплексных произведений,
отношений и углового управления. Они также сохраняют особенности
угловой координаты. Направления около границы диапазона фаз
остаются близкими на комплексной плоскости, хотя их числовые
значения могут заметно различаться. Этот вопрос рассмотрен
в подразделе~\ref{sec:ch4_complex_ops}; одно выделение угла
не делает произвольную обучаемую функцию периодической.
При нулевом модуле направление комплексного вектора физически
не определено.

При использовании относительных фаз все выбранные векторы
согласуются с общей опорой. Это уменьшает влияние общего
поворота фаз, сохраняя различия между электрическими величинами.
Сам выбор опорного напряжения и обработка его отсутствия
относятся к определению признакового пространства.
Они особенно существенны при провале напряжения и не заменяются
обычной нормализацией чисел.

\begin{table}[htbp]
\centering\small
\caption{Содержание рассматриваемых входных представлений}
\label{tab:ch4_input_representations}
\renewcommand{\arraystretch}{1.25}
\begin{tabularx}{\textwidth}{@{}>{\raggedright\arraybackslash}p{0.27\textwidth}>{\raggedright\arraybackslash}X@{}}
\toprule
Представление & Информация, передаваемая модели \\
\midrule
Мгновенные значения & Последовательность отсчётов согласованных каналов, включая форму переходного процесса \\
\addlinespace[4pt]
Декартовы спектральные координаты & Действительная и мнимая части выделенных комплексных коэффициентов \\
\addlinespace[4pt]
Полярные спектральные координаты & Модули и фазы фазных величин либо симметричных составляющих \\
\addlinespace[4pt]
Дополнительные преобразования & Произведения сигналов, ортогональные координаты и другие явно заданные производные признаки \\
\bottomrule
\end{tabularx}
\end{table}

(Примечание автора: Мгновенные значения – не понял, зачем так описывать про форму и т.д. Это прост набор точек.)

\subsection{Симметричные составляющие и производные признаки}
\label{sec:ch4_sequence_features}

На основе основной гармоники определяются векторы прямой,
обратной и нулевой последовательностей. Это переносит в признаки
известное описание несимметрии трёхфазного режима.
Нулевая последовательность связана с суммой фазных величин,
а разделение прямой и обратной последовательностей учитывает
их относительные фазы. Обратная последовательность может быть
существенна при несимметричных повреждениях, но её наличие
само по себе не является достаточным признаком конкретного события.

(Примечание автора: Да и можно бы добавить что-то из серии, что последовательности 1,2,0 - уже являются классикой в РЗА и используются постоянно в защитах по удобству физики понимания и т.д. Ну чёто красивое, да, понятно, что пока это лишь набросок раздела.)

Для произвольной трёхфазной величины обозначим комплексные
коэффициенты основной гармоники $\dot{X}_A$, $\dot{X}_B$
и $\dot{X}_C$. При $a=e^{j2\pi/3}$ составляющие определяются как
\begin{equation}
\begin{aligned}
 \dot{X}_0&=\tfrac13(\dot{X}_A+\dot{X}_B+\dot{X}_C),\\
 \dot{X}_1&=\tfrac13(\dot{X}_A+a\dot{X}_B+a^2\dot{X}_C),\\
 \dot{X}_2&=\tfrac13(\dot{X}_A+a^2\dot{X}_B+a\dot{X}_C).
\end{aligned}
\label{eq:ch4_sequence_features}
\end{equation}
Эти операции применяются отдельно к токам и напряжениям.
Полученные векторы также могут быть представлены декартовыми
или полярными координатами. Переход к симметричным составляющим
не определяет автоматически единственный лучший формат входа.

(Примечание автора: Про Кларка – надо бы чутка пояснить, а то я, например, вообще ничего не знал про этом, и мол нафига используем?)

Среди дополнительных представлений рассматриваются произведения
токов и напряжений и преобразование Кларка. Последнее формирует
ортогональные координаты трёхфазной системы. При несимметрии
и ОЗЗ отдельно учитывается нулевая составляющая; её смысл
не исчерпывается двумя координатами плоскости.
Произведения сигналов дают другой вид производных признаков,
физическая интерпретация которых зависит от того, используются
мгновенные значения или комплексные коэффициенты.

Предварительное вычисление таких величин изменяет вход
классификатора. В архитектуре PhysicsKAN часть операций
включена непосредственно в граф сети. Совпадение исходных
измерений ещё не делает эти варианты одинаковыми.
В первом случае преобразование задаётся подготовкой входа,
во втором включено в совместную обработку с обучаемыми слоями.
Общее сравнение поэтому сохраняет сведения и о составе признаков,
и о расположении операций.

\subsection{Масштабирование и преобразования при обучении}
\label{sec:ch4_feature_scaling}

Первый этап масштабирования согласует электрические величины
записей с разных объектов. Токи и напряжения приводятся
к выбранным базовым значениям с учётом коэффициентов
измерительных цепей. Это позволяет сопоставлять относительные
изменения режима, не смешивая их с различием номинальных уровней.
Коэффициенты для фазных и нулевых составляющих имеют собственный
физический смысл и не подменяются одним общим множителем.

Статистическая нормализация выполняет другую задачу.
Центрирование и приведение к выбранному масштабу подготавливают
распределения признаков для обучения. Их параметры определяются
по обучающей выборке и затем используются при оценке.
Если коэффициенты повторно вычислять по всей проверяемой записи,
во вход текущего момента может попасть информация о более поздней
части процесса. Поэтому способ определения масштаба является
частью протокола, а не только техническим свойством файла.

Для расширения разнообразия примеров в исследовании применяются
изменение амплитуд, общая инверсия полярности и циклическая
перестановка фаз. Совместная перестановка токов и напряжений
сохраняет соответствие фазных пар. Независимое перемешивание
этих каналов создало бы другое электрическое соотношение.
При изменении полярности также сохраняется согласованность
сигналов, из которых рассчитываются производные признаки.

Такие преобразования связаны с целевой меткой. Общее изменение
полярности согласованных токов и напряжений не равнозначно
инверсии только одного из них. Поэтому допустимое изменение
входа для классификации наличия события не переносится
автоматически на задачу определения направления РНМ.
Для каждого исследования состав преобразований рассматривается
вместе с его разметкой и применимостью измерений.

\subsection{Порядок сопоставления представлений}
\label{sec:ch4_feature_comparison}

В исследовании отдельно изменяются доступная предыстория,
набор сигналов и способ их преобразования.
Обработка мгновенных значений сравнивается со спектральными
координатами, фазные величины с симметричными составляющими,
а исходные признаки с дополненными физическими величинами.
При этом сохраняется различие семейства архитектуры,
конкретной конфигурации и способа формирования входа.

Оценка учитывает не только показатели распознавания,
но и длительность требуемого фрагмента, объём входных данных
и затраты на их вычисление. Для БАВР существенны также
доступность измерений и положение решения во времени.
Единые определения показателей даны
в разделе~\ref{sec:ch4_evaluation_methodology}.

\workinprogress{Количественные результаты сравнения предыстории и признаковых пространств, соответствующие таблицы и экспериментальные графики второй статьи ещё не включены. Текущая партия завершает описание подготовки и смысла входных представлений, а не выбор лучшего варианта по результатам опытов.}
